\chapter{Профилирование и сложность исходной программы}

\section{Профиль исходной программы}

Исходная программа (шаг~00, коммит 1f70c33) сжимает весь корпус на
компьютере debian за 10,10~с (28,9~МБ/с) и распаковывает за 2,07~с
(141~МБ/с). В~тактах процессора сжатие корпуса занимает 41,9~млрд тактов,
распаковка "--- 7,6~млрд тактов (таблица~\ref{tab:final-cycles}).

Профиль снимался на файле dickens (10,2~МБ английского текста). Основные
метрики счётчиков процессора, полученные командой \texttt{perf stat},
приведены в таблице~\ref{tab:base-perf}, а вывод команды "--- на
рисунке~\ref{fig:base-perfstat}.

\begin{table}[H]
\caption{Метрики исходной программы на файле dickens (perf stat, среднее
из 5 запусков)}
\label{tab:base-perf}
\begin{tabletext}
\centering
\begin{tabular}{|l|r|r|r|r|}
\hline
\multirow{2}{*}{\textbf{Метрика}} & \multicolumn{2}{c|}{\textbf{Сжатие}} & \multicolumn{2}{c|}{\textbf{Распаковка}} \\
\cline{2-5}
 & \textbf{debian} & \textbf{hppavillion} & \textbf{debian} & \textbf{hppavillion} \\
\hline
Такты, млн & 1150 & 1215 & 350 & 348 \\
Команды, млн & 2747 & 2639 & 862 & 811 \\
Команд за такт (IPC) & 2,39 & 2,17 & 2,46 & 2,33 \\
Ветвления, млн & 541 & 542 & 153 & "--- \\
Ошибки предсказания ветвлений, \% & 2,92 & 3,26 & 2,42 & "--- \\
Промахи кеша L1D, \% & 2,12 & "--- & 0,02 & "--- \\
Время, с & 0,279 & "--- & 0,092 & "--- \\
\hline
\end{tabular}
\end{tabletext}
\end{table}

\screenshot{base-perfstat}{Вывод perf stat для сжатия файла dickens исходной программой}{fig:base-perfstat}{%
Файл \texttt{bench/log/debian/00-baseline/raw/profile/perf-stat-compress.txt}:
\texttt{cat} в терминале. Или заново, на коммите 1f70c33 после
\texttt{make profile}: \texttt{perf stat -r 5 -e cycles:u,instructions:u,branches:u,branch-misses:u
bench/bin/prof --mode=compress --repeat=1 bench/corpus/silesia/dickens}.}

Распределение времени по функциям показывает плоский профиль gprof
(таблица~\ref{tab:base-gprof}, рисунок~\ref{fig:base-gprof}). Для него
программа собирается с ключами \texttt{-pg -fno-inline}, чтобы каждая
маленькая функция была видна отдельно; профиль охватывает 5~циклов сжатия и
распаковки.

\begin{table}[H]
\caption{Плоский профиль gprof исходной программы (debian, dickens)}
\label{tab:base-gprof}
\begin{tabletext}
\centering
\begin{tabular}{|l|r|r|l|}
\hline
\textbf{Функция} & \textbf{Доля времени, \%} & \textbf{Вызовов} & \textbf{Назначение} \\
\hline
\texttt{lz77\_compress} & 41,1 & 5 & цикл сжатия, поиск совпадений \\
\texttt{lz77\_decompress} & 16,5 & 5 & цикл распаковки \\
\texttt{bw\_flush} & 11,0 & 38\,841\,340 & вывод готовых байт архива \\
\texttt{br\_read} & 11,0 & 38\,841\,345 & чтение полей архива \\
\texttt{hash3} & 10,3 & 65\,453\,735 & хеш трёх байт \\
\texttt{ht\_insert} & 2,7 & 50\,962\,220 & вставка позиции в хеш-таблицу \\
\texttt{byte\_at} & 2,1 & 366\,666\,790 & байт окна или буфера \\
прочие & 5,3 & & \\
\hline
\end{tabular}
\end{tabletext}
\end{table}

\begin{figure}[H]
\centering
\begin{tikzpicture}
\begin{axis}[
  xbar, width=0.8\textwidth, height=6.5cm,
  xmin=0, xmax=50, xlabel={Доля времени, \%},
  symbolic y coords={прочие,byte\_at,ht\_insert,hash3,br\_read,bw\_flush,lz77\_decompress,lz77\_compress},
  ytick=data, yticklabel style={font=\small\ttfamily},
  nodes near coords, nodes near coords style={font=\small},
  bar width=9pt, enlarge y limits=0.08,
]
\addplot[fill=gray!60] coordinates {(41.1,lz77\_compress) (16.5,lz77\_decompress)
  (11.0,bw\_flush) (11.0,br\_read) (10.3,hash3) (2.7,ht\_insert)
  (2.1,byte\_at) (5.3,прочие)};
\end{axis}
\end{tikzpicture}
\caption{Распределение времени исходной программы по функциям (gprof)}
\label{fig:base-gprof}
\end{figure}

\screenshot{base-gprof}{Плоский профиль gprof исходной программы}{fig:base-gprof-shot}{%
\texttt{head -25 bench/log/debian/00-baseline/raw/profile/gprof.txt}
в терминале (строки <<Flat profile>> с первыми десятью функциями).}

Профиль по строкам исходного кода (\texttt{perf record}, рисунок~\ref{fig:base-perflines})
уточняет картину:
\begin{itemize}
\item 15--20~\% тактов приходится на побайтное сравнение во внутреннем
цикле поиска совпадения (строки 83--86 файла \texttt{lz77.c}), где каждый
байт кандидата читается функцией \texttt{byte\_at} с ветвлением и
взятием остатка от деления;
\item около 12~\% "--- на функцию \texttt{memmove} из glibc: после каждого
токена программа сдвигает упреждающий буфер влево;
\item заметную долю дают \texttt{fputc} и \texttt{fgetc}: каждый байт
архива и каждый байт распакованных данных проходит через отдельный вызов
библиотеки ввода-вывода.
\end{itemize}

\screenshot{base-perflines}{Самые <<горячие>> строки исходного кода (perf report)}{fig:base-perflines}{%
\texttt{head -30 bench/log/debian/00-baseline/raw/profile/perf-lines.txt},
или \texttt{perf report -i perf.data --sort srcline} после \texttt{make profile}.
Должны быть видны \texttt{lz77.c:84} и строки \texttt{memmove-vec-unaligned-erms.S}.}

Из числа вызовов функций в профиле gprof можно точно восстановить состав
потока токенов для файла dickens (таблица~\ref{tab:tokens}). Функция
\texttt{bw\_write} вызывается дважды на литерал и трижды на ссылку, а
\texttt{hash3} "--- один раз на каждую вставку в хеш-таблицу и один раз на
каждый токен. Рассчитанный по этим числам размер архива,
$(9\cdot926\,641 + 21\cdot1\,971\,662)/8 \approx 6\,218\,084$ байт, совпадает с
фактическим (6\,218\,082 байт) с точностью до округления.

\begin{table}[H]
\caption{Поток токенов при сжатии файла dickens (на один запуск)}
\label{tab:tokens}
\begin{tabletext}
\centering
\begin{tabular}{|l|r|}
\hline
\textbf{Величина} & \textbf{Значение} \\
\hline
Размер входного файла $n$, байт & 10\,192\,446 \\
Число токенов $T$ & 2\,898\,303 \\
\quad из них литералов & 926\,641 \\
\quad из них ссылок & 1\,971\,662 \\
Средняя длина ссылки, байт & 4,7 \\
Просмотрено кандидатов в цепочках, всего & 16\,170\,300 \\
Кандидатов на токен, в среднем & 5,6 \\
Прочитано байт при сравнении, на кандидата & 4,5 \\
\hline
\end{tabular}
\end{tabletext}
\end{table}

Наконец, чтобы отделить качество кода от настроек компилятора, исходная
программа была собрана с разными уровнями оптимизации GCC
(таблица~\ref{tab:base-variants}, \texttt{make bench-compare}). Переход от
\texttt{-O0} к \texttt{-O1} ускоряет сжатие вдвое, но дальнейшие уровни
дают лишь несколько процентов. Значит, оставшееся время определяется самим
кодом "--- лишними копированиями и вызовами, "--- и его нужно менять.

\begin{table}[H]
\caption{Исходная программа с разными флагами компилятора (debian, весь корпус)}
\label{tab:base-variants}
\begin{tabletext}
\centering
\begin{tabular}{|l|r|r|r|r|}
\hline
\textbf{Флаги} & \textbf{Сжатие, с} & \textbf{Распаковка, с} & \textbf{Сжатие, МБ/с} & \textbf{Распаковка, МБ/с} \\
\hline
\texttt{-O0} & 21,88 & 3,13 & 13,3 & 93,3 \\
\texttt{-O1} & 10,90 & 2,07 & 26,8 & 141,0 \\
\texttt{-O2} & 10,47 & 2,07 & 27,9 & 140,9 \\
\texttt{-O3} & 10,12 & 1,99 & 28,9 & 146,8 \\
\texttt{-O3 -march=native} & 10,62 & 2,02 & 27,5 & 144,4 \\
\quad + LTO & 9,50 & 1,76 & 30,7 & 166,3 \\
\quad + LTO + PGO & 9,18 & 1,75 & 31,8 & 166,8 \\
\hline
\end{tabular}
\end{tabletext}
\end{table}

\section{Оценка алгоритмической сложности исходной программы}
\label{sec:base-complexity}

Сложность оценивается в асимптотических обозначениях
как число элементарных операций в зависимости от размера входа.
Обозначим $n$ "--- размер входного файла, $T$ "--- число токенов
($n/M \le T \le n$), остальные обозначения "--- по
таблице~\ref{tab:params}.

\textbf{Сжатие.} Для каждого токена программа выполняет:
\begin{enumerate}
\item вычисление хеша и выбор головы цепочки "--- $O(1)$;
\item просмотр не более $C$ кандидатов цепочки, для каждого "--- сравнение
не более $M$ байт, каждый через \texttt{byte\_at} "--- $O(C \cdot M)$;
\item запись токена в архив "--- $O(1)$;
\item вставку в хеш-таблицу каждой из пройденных позиций "--- $O(1)$ на байт
входа;
\item сдвиг упреждающего буфера функцией \texttt{memmove} на длину
токена: копируется весь остаток буфера, от 258 до $L = 4096$ байт, "---
$O(L)$.
\end{enumerate}

Суммарное время сжатия
\begin{equation}
\label{eq:base-compress}
T_{\text{сж}}(n) = O\bigl(n + T \cdot (C \cdot M + L)\bigr) =
O\bigl(n \cdot (C \cdot M + L)\bigr),
\end{equation}
где худший случай $T = n$ достигается на несжимаемых данных (одни
литералы). Поскольку $C$, $M$ и $L$ "--- константы алгоритма, асимптотически
сжатие линейно: $T_{\text{сж}}(n) = O(n)$. Однако константа очень велика:
в худшем случае до $C\cdot M + L \approx 70\,000$ операций на байт.

На реальных данных цепочки короткие: для dickens в среднем просматривается
5,6~кандидата на токен, и сравнение каждого обрывается через 4,5~байта
(таблица~\ref{tab:tokens}), то есть поиск стоит около 25 сравнений на токен.
Сдвиг буфера стоит намного дороже: размер упреждающего буфера равномерно
убывает от 4096 до 258 байт и затем снова дополняется, так что в среднем
на каждый токен копируется около $(4096 + 258)/2 \approx 2200$ байт. Для
dickens это примерно $2{,}9\cdot10^6 \cdot 2200 \approx 6{,}3$~ГБ копирования
на 10~МБ входных данных "--- более 600 байт на каждый входной байт.
Именно этот член $O(T\cdot L)$ определяет время на данных с короткими
совпадениями.

\textbf{Распаковка.} Каждый токен декодируется за $O(1)$, каждый выходной
байт записывается за $O(1)$: $T_{\text{расп}}(n) = O(n)$. Константу
определяют вызовы \texttt{fgetc} и \texttt{fputc} "--- по одному на каждый
байт архива и каждый байт результата.

\textbf{Память.} Сжатие использует таблицу голов цепочек ($H$ элементов
по 8~байт, 512~КиБ), массив \texttt{ht\_next} ($W$ элементов, 32~КиБ), окно
и упреждающий буфер (по 4~КиБ): $O(H + W + L) = O(1)$ "--- память не
зависит от размера файла. Распаковка хранит только окно, $O(W) = O(1)$.

\textbf{Экспериментальная проверка.} Линейность подтверждается
измерением: префиксы синтетических файлов размером от 1 до 20~МБ сжимались
исходной программой, и для каждого считалось число выполненных команд
процессора (таблица~\ref{tab:ipb} в~главе~6). Число команд на байт входа
постоянно с точностью до 1~\%: около 251 для текста, 384 для случайных
данных и 33 для нулевых байт. Случайные данные "--- худший случай: в них
почти нет совпадений, каждый байт становится литералом и вызывает сдвиг
буфера.
